\chapter{Getting a structure}
\label{ch:structure}

Every preparation starts with a three-dimensional structure of your protein.
ProPrep can download one from the Protein Data Bank, fetch a model predicted
by AlphaFold, or open a file that is already on your computer.
This chapter shows each route. It also explains the difference between the
``asymmetric unit'' and the ``biological assembly'', which decides whether you
get the whole working molecule or only part of it.

\begin{glance}
\glanceitem{What it does}{Downloads or opens a protein structure, shows its
  description, and, if needed, rebuilds the complete working molecule from
  crystal symmetry.}
\glanceitem{You need it when}{Always. Every other tool works on the structure
  you load here.}
\glanceitem{Before you start}{An internet connection for downloads. A new,
  empty project folder for each protein (Chapter~\ref{ch:projectfolder}).}
\glanceitem{Where to find it}{\menupath{1a}{Structure Loader} and
  \menupath{1b}{Biological Assembly Generator}. In Workflow Mode, loading is
  the first step of every workflow.}
\glanceitem{You get}{The downloaded file, for example \file{4HHB.pdb}.
  Sometimes also a copy without hydrogens (\file{9ZZZ_noH.pdb}), a rebuilt
  assembly from 1b (\file{1HHO_biological_assembly.pdb}) or an AlphaFold model
  (\file{AF-P69905-F1-model_v6.cif}).}
\end{glance}

\section{Two things to know first}
\label{sec:structure-basics}

\subsection{Asymmetric unit or biological assembly?}
\label{sec:structure-asu}

This is the point that confuses most beginners, so it is worth two minutes.

A protein crystal is a regular, repeating stack of molecules. Crystallographers
do not deposit the whole crystal. They deposit the smallest piece from which the
whole crystal can be rebuilt by symmetry. This piece is the
\term{asymmetric unit}: the set of chains that are actually written in the PDB
file. The asymmetric unit is a property of the crystal, not of the protein. It
may contain exactly one working molecule, only part of one, or several copies.

The \term{biological assembly} is the molecule that is believed to exist and
work in the cell or in solution. For hemoglobin this is the
$\alpha_2\beta_2$ tetramer: four chains, two alpha and two beta. The PDB file
contains instructions, in its \term{REMARK 350} lines, for building the
biological assembly from the asymmetric unit. These instructions are rotations
and shifts that place copies of the deposited chains.

Two hemoglobin entries show the difference:

\begin{tabularx}{\linewidth}{@{}lLL@{}}
\toprule
\textbf{Entry} & \textbf{Asymmetric unit (what the file holds)} & \textbf{Biological assembly} \\
\midrule
\pdb{4HHB} & All four chains A, B, C, D. & The same four chains. Nothing to build. \\
\pdb{1HHO} & One alpha and one beta chain (A, B) only: half the molecule. & Four chains: the deposited pair plus a symmetry copy. \\
\bottomrule
\end{tabularx}

\medskip
So for \pdb{1HHO} the downloaded file is \emph{half} a hemoglobin. If you
simulate it as it is, you simulate something that does not exist in nature. The
reverse also happens: an asymmetric unit can hold two or more copies of a
protein that works alone, packed together only by the crystal. Then you keep one
copy later, in the PDB Filter (Chapter~\ref{ch:cleaning}).

Structures from NMR and from AlphaFold are not crystals. For them the question
does not arise, but note that an AlphaFold model is always a \emph{single
chain}: the AlphaFold model of hemoglobin alpha (\pdb{P69905}) is one alpha
chain, not the tetramer.

\begin{tip}
How do you know which you need? Read the entry's page on the RCSB website, or
look at ProPrep's own information: the metadata summary
(Section~\ref{sec:structure-metadata}) says whether the file contains
\prompt{biological assembly information (REMARK 350)}, and tool \key{1b} names
the assembly, for example \prompt{TETRAMERIC (author determined)}. If that
assembly has more chains than the file, you need the assembly.
\end{tip}

\subsection{PDB format and mmCIF format}

Structures come in two file formats. The older \term{PDB format} (files ending
in \file{.pdb}) is what most of ProPrep expects. The newer \term{mmCIF} format
(files ending in \file{.cif}) is used by the AlphaFold and AlphaFill downloads
and by very large entries. In version 1.23.0 the file browsers list and accept
only \file{.pdb} files. Wherever you have a choice, use PDB format.

\section{Step by step: download a structure by its PDB ID}
\label{sec:structure-pdbid}

This is the route most people use. The example is human deoxyhemoglobin,
\pdb{4HHB}.

\begin{steps}
  \item In the main menu, type \key{1a} (Structure Loader). The menu shows four
        options; only the first is needed now.
\begin{screen}
══ Structure Loader Menu ══
  1. ✓ Load protein structure
  2. ✓ Load AMBER topology & coordinate files
  3. ○ View structure metadata
       ⚠ Load a structure first
  4. ✓ Load a ligand file or library (SDF/mol2/SMILES) for docking
...
Enter your choice (1): §ans«1»
\end{screen}
  \item ProPrep asks where the structure should come from. Type \key{1}
        (Protein Data Bank), then \key{1} again (Enter PDB ID directly).
\begin{screen}
═══ Load Structure - Select Source ═══

1. Query Protein Data Bank (RCSB)
2. Query AlphaFold Database
3. Query AlphaFill Database (AlphaFold + ligands/cofactors)
4. Load local PDB file
5. View local files together (spatial comparison, view-only)
6. Done (return to the Structure Loader menu)

Select source (1): §ans«1»

═══ Query Protein Data Bank (RCSB) ═══

1. Enter PDB ID directly
2. Search by entry title keyword
3. Search by gene/protein (UniProt)
4. Search by ligand (code or name)
5. Back to source selection

Select method (1): §ans«1»
\end{screen}
  \item Type the four-character code, \typed{4hhb}. Lower case is fine;
        ProPrep changes it to upper case.
  \item Choose the structure type. Type \key{1} for the asymmetric unit, as
        deposited, or \key{2} for the biological assembly
        (Section~\ref{sec:structure-asu}). For \pdb{4HHB} the two are the same,
        so take \key{1}.
\begin{screen}
Enter PDB ID (e.g., 1ABC, 7XYZ): §ans«4hhb»

Structure type:
  1. Asymmetric unit (default - as deposited)
  2. Biological assembly (functional oligomeric form)
Select structure type (1): §ans«1»

Downloading 4HHB...

✓ Downloaded 4HHB

PDB ID: 4HHB
Classification: OXYGEN TRANSPORT
Deposition Date: 07-MAR-84

Title: THE CRYSTAL STRUCTURE OF HUMAN DEOXYHAEMOGLOBIN AT 1.74 ANGSTROMS
RESOLUTION

Experimental Method: X-RAY DIFFRACTION
Resolution: 1.74 Å

Publication:
  THE CRYSTAL STRUCTURE OF HUMAN DEOXYHAEMOGLOBIN AT 1.74 A RESOLUTION
  G.FERMI,M.F.PERUTZ,B.SHAANAN,R.FOURME
  J.MOL.BIOL. V. 175 159 1984
  PubMed: https://pubmed.ncbi.nlm.nih.gov/6726807/
  DOI: https://doi.org/10.1016/0022-2836(84)90472-8
  Abstract and free full text, if any: View structure metadata
\end{screen}
        The download takes about two seconds. The short summary tells you the
        method and the \term{resolution} (how fine the detail of an X-ray or
        cryo-EM structure is; smaller numbers are better, and 2~Å or less is
        good).
  \item ProPrep returns to the source menu, so that you can load more. Type
        \key{6} (Done) to go back to the Structure Loader menu.
\end{steps}

The file \file{4HHB.pdb} is now in your project folder, exactly as the Protein
Data Bank provides it. Notice that ProPrep did not ask which chains, ligands or
waters you want. You choose those later, in the PDB Filter
(Chapter~\ref{ch:cleaning}).

\subsection{Two automatic checks after a download}

Depending on the entry, ProPrep may stop with one or two extra screens.

\paragraph{Residue numbering warning.} When a chain does not start at residue
number~1, a box lists those chains, for example
\prompt{Chain A: First residue is LYS \#13}. This is information, not an error.
The usual reasons, which the box also names, are missing residues at the start
of the chain, numbering that counts a signal (leader) sequence, or numbering
taken from a reference protein. Keep it in mind when you later type residue
numbers.

\paragraph{Hydrogen removal.} Most crystal structures have no hydrogen atoms.
NMR structures, very-high-resolution crystal structures and many cryo-EM
entries do. Hydrogen names in PDB files often do not match what the
\term{force field} (the set of parameters that describes the atoms in a
simulation) expects, so ProPrep offers to remove them. The program that builds
the simulation, \term{tLEaP}, adds them back later with the right names. This
screen came from the cryo-EM entry \pdb{9ZZZ}:
\begin{screen}
Hydrogen Atom Removal
This structure contains 14334 hydrogen atoms. Hydrogen naming in PDB structures
is often inconsistent with force field expectations. It is recommended to remove
them and let tLEaP add them back with correct naming, positions, and protonation
states.

Remove hydrogen atoms? [y/n] (y): §ans«y»
H-stripped structure saved to: 9ZZZ_noH.pdb
Removed 14334 hydrogen atoms
\end{screen}
Answer \key{y} in most cases. ProPrep keeps the original file and writes a
second file, \file{9ZZZ_noH.pdb}, which contains only the coordinates (no
header lines).

\begin{watchout}
Three things to know about hydrogen removal in version 1.23.0:
\begin{itemize}
  \item Your answer is remembered for the rest of the session. Later downloads
        with hydrogens are stripped without asking; the screen then says
        \prompt{(using previous selection)}.
  \item The \file{_noH} file outranks every structure you load afterwards
        (Section~\ref{sec:structure-switching}). It also has no REMARK 350 lines,
        so the Biological Assembly Generator cannot use it.
  \item On the alternate-location entry \pdb{3NIR}, 18 hydrogen atoms (on one
        conformation of Pro22) stayed in \file{3NIR_noH.pdb}.
\end{itemize}
\end{watchout}

\section{Getting the complete working molecule}
\label{sec:structure-assembly}

When the asymmetric unit is not the working molecule, you have two ways to get
the biological assembly. Both were tried with \pdb{1HHO}, whose file holds only
one alpha and one beta chain.

\subsection{Route 1 (recommended for simulations): the Biological Assembly Generator}
\label{sec:structure-1b}

Tool \key{1b} reads the REMARK 350 instructions in the downloaded file and
builds the assembly in one model with distinct chain letters.

\begin{steps}
  \item Start a new project. Load \pdb{1HHO} with \key{1a} \textrightarrow{}
        \key{1} \textrightarrow{} \key{1} \textrightarrow{} \key{1}, typing
        \typed{1HHO} and choosing \key{1} (asymmetric unit). Type \key{6} and
        then \key{m} to return to the main menu.
  \item Type \key{1b}, then \key{1}
        (\prompt{Generate biological assembly from asymmetric unit}).
  \item Check the first line: it names the file ProPrep is using. It must be
        your download (here \file{1HHO.pdb}).
\begin{screen}
Using RCSB PDB PDB file: .../asm/1HHO.pdb

Input structure: 1HHO.pdb
Source: rcsb_pdb_file

Found 1 biological assembly definition(s) in REMARK 350

Available biological assemblies:
  1. TETRAMERIC (author determined) (software determined) -
     2 transform(s), chains: A, B

Options:
  1-N  Select assembly from REMARK 350
  m    Enter manual transformation matrix
  c    Cancel

Select option (1): §ans«1»
Include ligands/cofactors (HETATM)? [y/n] (y):
Include water molecules? [y/n] (n):
\end{screen}
        The line \prompt{TETRAMERIC \ldots{} 2 transform(s), chains: A, B} is
        the \term{quaternary structure} description: the assembly is a tetramer,
        built by applying two transformations to chains A and B.
  \item Type \key{1} to choose that assembly. Press \key{Enter} twice to keep
        ligands and cofactors (here the four hemes) and to leave out the
        crystal waters.
\begin{screen}
Generating biological assembly...

Assembly generated successfully!
  Original: 2 chains, 2,287 atoms
  Assembly: 4 chains (A, B, C, D), 4,574 atoms
  Transforms applied: 2

Saved: .../asm/1HHO_biological_assembly.pdb
Registered in workspace as 'biological_assembly_pdb_file'
\end{screen}
\end{steps}

It takes less than a second. The new file \file{1HHO_biological_assembly.pdb}
holds one model with chains A, B, C and D, four hemes, and no water. From now
on, later tools use this file automatically.

\begin{watchout}
Run \key{1b} only once per project. When you open it again, it takes the
assembly it has made as its input and offers the same
\prompt{TETRAMERIC} assembly. Accepting would apply the symmetry a second time
to a structure that is already complete. Type \key{c} (Cancel) instead.
\end{watchout}

If you type a number that does not exist, for example \key{7}, ProPrep prints
\prompt{Invalid assembly number: 7} and goes back to the menu without asking
again. If an assembly would need more than 62 chains, ProPrep warns that the PDB
format has only 62 chain letters and asks whether to continue (this was not
tried for this manual).

\begin{watchout}
In version 1.23.0, choosing \key{m} (\prompt{Enter manual transformation
matrix}) closes ProPrep completely. After you type the chains, the program
stops with:
\begin{screen}
Enter transformation matrix (3x3 rotation matrix):
Error executing Show Generate biological assembly from asymmetric unit menu:
name 'self' is not defined
...
Error: name 'self' is not defined
Run with --debug flag for more details
[Session recording stopped]
\end{screen}
Files already written stay on disk. Do not use \key{m}. If \key{1b} offers only
\key{m} and \key{c} because it says \prompt{No REMARK 350 records found in PDB
file.}, it is reading a file without header lines, usually a \file{_noH} file
(Section~\ref{sec:structure-switching}). Start a new project and download the
assembly directly (Route~2).
\end{watchout}

\subsection{Route 2: download the assembly directly}

At step~4 of Section~\ref{sec:structure-pdbid}, type \key{2}
(Biological assembly) instead of \key{1}:
\begin{screen}
Select structure type (1): §ans«2»

Checking available biological assemblies for 1HHO...
Found 1 biological assembly

Downloading biological assembly 1...

✓ Downloaded biological assembly 1 for 1HHO
...
Assembly contains 4 chains, 4,792 atoms
\end{screen}
If an entry has several assemblies, ProPrep lists them and asks
\prompt{Select assembly number (1)}. If it has none, it says so and falls back
to the asymmetric unit. The file is called \file{1HHO_assembly1.pdb}.

This file comes straight from the RCSB, and it is built differently from the
1b file: it stores the tetramer as \emph{two models}, each with chains A and B,
and it keeps the crystal waters (221 of them). The summary also shows no
resolution line. Tools that read only the first model would see only two
chains. How the later ProPrep tools treat such a two-model file was not tested,
so for simulations prefer Route~1.

\section{Other ways to find and load a structure}
\label{sec:structure-other}

\subsection{Searching the Protein Data Bank by name}
\label{sec:structure-search}

Use this when you do not know the PDB code. At \prompt{Select method}, type
\key{2} (Search by entry title keyword), type a few words, and optionally an
organism. The example searched for horse cytochrome~c:
\begin{screen}
Enter search keyword (e.g., "carbonic anhydrase", "SOD1", "hemoglobin"):
§ans«cytochrome c»
Filter by organism? (optional, press Enter to skip) (): §ans«Equus caballus»

Searching PDB for "cytochrome c" in "Equus caballus"...
35 structures match
Maximum number of structures to list, most relevant first (range: 1-10000,
default: 100) (100):

═══ Search Results (Page 1/2) ═══
 #   PDB ID  Title                                          Method  Res (Å)  Metals
 1   6K9J    0.98 A three-dimensional structure of horse    X-ray   0.98     HEC (Fe)
...
 8   1HRC    HIGH-RESOLUTION THREE-DIMENSIONAL STRUCTURE    X-ray   1.90     HEC (Fe)
...
Options:
  • Select: single (1), multiple (1,3,5), range (1-3), or 'all'
  • Type 'next' to see next page
  • Type 'filter' to refine search
  • Type 'back' to return to the Structure Loader menu

Selection:
\end{screen}
(The table is shortened here.) Results are sorted best resolution first, with
NMR entries (resolution \prompt{N/A}) last. The \prompt{Metals} column lists
metal-containing groups, for example \prompt{HEC (Fe)} for the iron-containing
heme of cytochrome~c.

At \prompt{Selection:} you can type:
\begin{itemize}
  \item a number such as \typed{8} to download one entry. ProPrep shows the
        entry and asks \prompt{Download 1HRC? [y/n] (y)}. This route always
        downloads the asymmetric unit;
  \item several numbers (\typed{3,9}) or a range (\typed{1-3}) to download them
        together, after \prompt{Download all 2 structures? [y/n] (y)};
  \item \typed{next} and \typed{prev} to change page. The numbers start again
        at 1 on each page and refer to that page only;
  \item \typed{filter} to narrow the list: \key{1} maximum resolution,
        \key{2} method (X-ray, NMR, cryo-EM), \key{3} must contain a metal
        (for example \typed{Zn}). With the zinc filter the 35 horse
        cytochrome~c entries became 4;
  \item \typed{back} to return to the source menu.
\end{itemize}

\begin{tip}
Always give an organism. Several words without one search very broadly:
\typed{horse heart cytochrome c} alone matched 10\,104 entries.
\end{tip}

\begin{watchout}
Downloading several entries at once skips the summary, the numbering check and
the hydrogen question. The NMR entry \pdb{1AKK} kept its 878 hydrogens this way.
Only the \emph{last} entry of the batch becomes the structure that later tools
use. Download the entry you want to prepare on its own.
\end{watchout}

Two smaller points: \prompt{4. Clear all filters} only says
\prompt{To clear filters, start a new search}, and \typed{back} says
\prompt{Back to the Structure Loader menu} although it returns to the source
menu.

\subsection{Searching by gene or protein name (UniProt)}

\term{UniProt} is the main database of protein sequences; each protein has an
\term{accession} code such as \pdb{P00441}. At \prompt{Select method}, type
\key{3}, then a gene name, protein name or accession, and optionally an
organism. With \typed{SOD1} and \typed{human}, ProPrep showed the UniProt entry
(superoxide dismutase [Cu-Zn], 154 residues, cofactors
\prompt{Cu cation, Zn(2+)}), offered more detail, and then said
\prompt{This entry has 156 associated PDB structures}. A table of those
structures follows, 20 per page, with method, resolution and the chains and
residues each one covers (for example \prompt{A/F=2-154}). Typing a number
downloads that entry \emph{at once}, with no confirmation and no choice of
assembly.

\subsection{Searching by ligand}

To find structures that contain a particular small molecule, type \key{4} at
\prompt{Select method}. Every ligand in the PDB has a short code, for example
\prompt{HEM} for heme. You can type the code or words from its name. Names
match whole words only. For example, \typed{biopterin} does not find
the ligand named \typed{tetrahydrobiopterin}. ProPrep lists matching codes, can show related
codes (other protonation or oxidation states of the same compound), and then
asks which codes to search for and how many results to list. The results table
has two extra columns: \prompt{Ligand}, and \prompt{RSCC}, which tells you how
well the ligand fits the experimental electron density (1 is a perfect fit).

\subsection{Predicted structures from AlphaFold}
\label{sec:structure-alphafold}

\term{AlphaFold} predicts a protein's structure from its sequence. Use it when
there is no experimental structure, or when the experimental one is missing
large pieces. At \prompt{Select source}, type \key{2}, then \key{1}
(Retrieve by UniProt ID) and the accession:
\begin{screen}
Enter UniProt ID (e.g., P00520, Q5VSL9): §ans«p69905»

Checking AlphaFold Database for P69905...

✓ Retrieved AF-P69905-F1
  Protein: Hemoglobin subunit alpha
  Organism: Homo sapiens
  File: AF-P69905-F1-model_v6.cif
  Confidence: 98.0 (Very High)

View AlphaFold metadata? [y/n] (y): §ans«y»
...
Model Quality:
  Mean pLDDT: 98.0
  Quality: Very High

Confidence Distribution:
  Very High (>90): 1069 residues (99.3%)
  Confident (70-90): 0 residues (0.0%)
  Low (50-70): 8 residues (0.7%)
  Very Low (<50): 0 residues (0.0%)
\end{screen}

AlphaFold gives every residue a confidence score, the \term{pLDDT}, from 0 to
100. Above 90 the prediction is usually as reliable as a good experimental
structure. From 70 to 90 the backbone is generally right. Below 70, and
especially below 50, the region is uncertain and is often a flexible, disordered
part of the protein. The metadata screen also shows the sequence colored by
confidence (the colors appear on screen only). Look at the low-confidence
regions before you simulate them.

Option \key{2} (Search by gene name) lists matching UniProt entries, for example
\typed{HBB} gave ten beta-globins from different species, and you pick one by
number.

\begin{watchout}
Points to know about AlphaFold downloads in version 1.23.0:
\begin{itemize}
  \item The \prompt{Length} line and the confidence distribution count
        \emph{atoms}, not residues. Hemoglobin alpha has 142 residues, but the
        screen says \prompt{Length: 1077 residues}. The percentages are
        therefore per atom. After a gene-name search, the metadata shows
        \prompt{UniProt ID: N/A} and \prompt{Organism: Unknown}.
  \item The file is in mmCIF format (\file{.cif}). The file browsers of
        \key{1a} and \key{3b} do not list it and refuse it if you type its path.
        If you need the model in a browser (for example to align it,
        Chapter~\ref{ch:compare}), download the PDB-format file from the
        protein's page on the AlphaFold Database website and put it in your
        project folder.
  \item ProPrep keeps one AlphaFold model at a time; a second AlphaFold
        download replaces the first.
  \item An AlphaFold model contains one chain and no ligands, metals or
        waters.
  \item Without internet the message is \prompt{No AlphaFold structure found
        for P69905} and \prompt{Not all proteins have AlphaFold predictions
        yet}. The real cause may be the connection.
\end{itemize}
\end{watchout}

\subsection{AlphaFold models with ligands: AlphaFill}

\term{AlphaFill} takes an AlphaFold model and copies in (``transplants'')
ligands, cofactors and metal ions from similar experimental structures. At
\prompt{Select source}, type \key{3}; the questions are the same as for
AlphaFold. The file is saved as \file{alphafill_P69905.cif}.

\begin{watchout}
In version 1.23.0 the AlphaFill screen always says \prompt{No transplanted
ligands/cofactors found} and \prompt{This AlphaFill structure does not have
enrichment data}. Do not believe it. The file for \pdb{P69905} does contain
transplanted groups, including heme (HEM), carbon monoxide and oxygen. The
\prompt{Length: 168 residues} counts the 142 amino acids plus these groups. The
file itself is fine; look at it in the viewer (Chapter~\ref{ch:compare}).
\end{watchout}

\subsection{Your own file}
\label{sec:structure-local}

At \prompt{Select source}, type \key{4}. A simple file browser opens in your
project folder:
\begin{screen}
Load Local PDB File

Current directory: .../s3load
  [ 1] .. (parent)
  [ 2] 1AKK.pdb  (160.2 KB)
  [ 3] 1HHO_assembly1.pdb  (386.7 KB)
...
Commands:
  N          - open dir N, or select file N
  ..         - go up one level
  <path>     - jump to a file or directory path
  q          - cancel

Selection (q):
\end{screen}
Type the number of a file, or of a folder to open it, or \typed{..} to go up
one level. You can also type the full path of a folder or a file, for example
\typed{/Users/me/structures/1CRN.pdb}. A bare file name such as
\typed{1CRN.pdb} is not accepted (\prompt{Enter a single number, '..', or
'q'.}); type the full path or start it with \typed{./}.
\begin{screen}
✓ Loaded structure from local file
  File: 1CRN.pdb
\end{screen}
ProPrep \emph{copies} the file into your project folder and works on the copy;
your original stays untouched. If the file contains hydrogens, you get the
hydrogen question. If its name contains \file{AF-} or \file{alphafold}, ProPrep
treats it as an AlphaFold model.

The browser lists only \file{.pdb} files. Typing the path of another file gives
\prompt{junk.txt is not a PDB file.} or \prompt{1UBQ.cif is not a PDB file.}.
\key{q} opens a plain question, \prompt{Enter path to PDB file:}. That route
accepts a \file{.cif} file, but only half-loads it (\prompt{BioPython parsing
skipped}); it was not checked which later tools then work. Convert your file
to PDB format first.

\subsection{Looking at several files together}

Source \key{5} (View local files together) opens two or more local PDB files in
the 3D viewer, in their real coordinates. You choose files in the same browser,
several at once with \typed{2,5} or \typed{1-5}. This is for looking only:
nothing is loaded for preparation, and the files are not superimposed, so
unrelated structures may appear far apart. To superimpose structures, use the
Structure Aligner (Chapter~\ref{ch:compare}).

\subsection{Starting ProPrep with a structure}

You can download or load a structure when you start ProPrep:
\begin{terminal}
proprep --full-menu --pdbid 1CRN
proprep --full-menu --pdbfile my_protein.pdb
\end{terminal}
ProPrep then asks no project name and works in the folder where you started it.
If the download fails (for example with no internet), the screen says
\prompt{Failed to download PDB 1CRN}, followed by a misleading log line
\prompt{Downloaded and loaded PDB: 1CRN}; the main menu then shows nothing
loaded. Chapter~\ref{ch:cmdline} lists these options.

\subsection{Other Structure Loader options}

Option \key{2} (Load AMBER topology \& coordinate files) loads the files of a
finished simulation system, for viewing a trajectory
(Chapters~\ref{ch:analysis} and~\ref{ch:compare}). Option \key{4} loads a ligand
file for docking (Chapter~\ref{ch:docking}).

\section{Reading the description of a structure}
\label{sec:structure-metadata}

Before you prepare a structure, check what is in it. In the Structure Loader
menu, type \key{3} (View structure metadata). If several kinds of structure are
loaded, ProPrep asks which one (\prompt{1. RCSB PDB structure}, \prompt{2. Local
PDB structure}, and so on). Opening the metadata also opens the 3D viewer.
\begin{screen}
===== PDB METADATA VIEWER =====
File: .../s3load/1MFM.pdb
ID: 1MFM

Available Sections:
 1. Summary                  - Quick overview of structure and experimental details
 2. Safety & Status Records  - Obsolete/superseded status and warnings
 3. Header Records           - Classification, deposition date, and keywords
...
 8. Hetero Compounds         - Ligands, cofactors, and non-standard residues
...
13. Export Metadata          - Save metadata to JSON or text file
14. Help                     - Detailed explanations of what each section contains
15. Exit                     - Return to previous menu
\end{screen}

The most useful sections are:
\begin{description}
  \item[1 Summary] The \term{composition} of the entry in the depositors' words.
        For the superoxide dismutase \pdb{1MFM} it read
        \prompt{Compound: MOL\_ID: 1; MOLECULE: PROTEIN (COPPER,ZINC
        SUPEROXIDE DISMUTASE); CHAIN: A; \ldots{} MUTATION: YES}, followed by the
        source organism and expression system, the number of chains and
        residues (\prompt{Chain A: 153 residues}), and which REMARK records
        exist, including \prompt{biological assembly information (REMARK 350)}.
        Read the \prompt{MUTATION: YES} and \prompt{ENGINEERED: YES} flags: the
        deposited protein may differ from the natural one.
  \item[5 Sequence Information] The sequence of each chain.
  \item[7 Citations] The publication, with PubMed and DOI links.
  \item[8 Hetero Compounds] Every non-protein group, by code, with name,
        formula and location, for example \prompt{CU}, \prompt{ZN}, and
        \prompt{HOH} (water). Use this list to decide what to keep in the PDB
        Filter.
  \item[13 Export Metadata] Saves everything to \file{1MFM_metadata.json} or a
        text file in your project folder.
\end{description}

For AlphaFold and AlphaFill structures the metadata screen is the one shown in
Section~\ref{sec:structure-alphafold}.

\section{Loading several structures and switching between them}
\label{sec:structure-switching}

ProPrep can hold one structure of each kind at the same time: one from the
RCSB, one local file, one AlphaFold model and one AlphaFill model. Loading
another of the same kind replaces the previous one. The Structure Loader menu
lists what is loaded, for example
\prompt{[All sources loaded: RCSB, Local, AlphaFold, AlphaFill]}.

Which of these does the next tool use? Not necessarily the one you loaded last.
Each tool picks its input by a fixed order of preference, and files that
ProPrep made itself come first: the hydrogen-stripped \file{_noH} file, an
assembly from \key{1b}, aligned files, and the filtered, repaired and
protonated files of Chapters~\ref{ch:cleaning} and~\ref{ch:protonation}. A fresh
download comes after all of them.

\begin{watchout}
\textbf{Loading a new structure does not make ProPrep use it.} In our tests:
\begin{itemize}
  \item After \pdb{9ZZZ} had been loaded with hydrogen removal, we loaded
        \pdb{4HHB}, \pdb{1HHO} and four more entries. Tool \key{1b} still said
        \prompt{Using H-Stripped PDB file: .../9ZZZ\_noH.pdb}.
  \item After building the \pdb{1HHO} assembly, we downloaded \pdb{4HHB}.
        Tool \key{1b} still used \file{1HHO_biological_assembly.pdb}.
\end{itemize}
Every tool prints a \prompt{Using \ldots{} file:} line at the top. Read it each
time. Some tools (Homology Searcher, Structure Viewer, PDB Filter) show a table
of the structures in the workspace and let you choose; \key{1b} does not.
\end{watchout}

The reliable rule is: \textbf{one protein, one project.} Start a new project
(a new folder) for each structure you prepare (Chapter~\ref{ch:projectfolder}).
The Workspace Options menu (\key{w}) has a \prompt{Reset workspace} option,
which should also clear the old files from ProPrep's memory; it was not tried
for this manual.

\section{Options at a glance}

\begin{longtable}{@{}p{0.33\linewidth}p{0.62\linewidth}@{}}
\toprule
\textbf{Question} & \textbf{Answers} \\
\midrule
\endhead
\prompt{Select source} & \key{1} RCSB, \key{2} AlphaFold, \key{3} AlphaFill,
  \key{4} local file, \key{5} view local files together, \key{6} done \\
\prompt{Select method} (RCSB) & \key{1} PDB ID, \key{2} title keyword,
  \key{3} gene or protein (UniProt), \key{4} ligand, \key{5} back \\
\prompt{Select structure type} & \key{1} asymmetric unit (default),
  \key{2} biological assembly from RCSB \\
\prompt{Remove hydrogen atoms?} & \key{y} (default) writes \file{<name>_noH.pdb};
  the answer is reused for the session \\
\prompt{Selection:} (search results) & \typed{3}, \typed{1,3,5}, \typed{1-3},
  \typed{all}, \typed{next}, \typed{prev}, \typed{filter}, \typed{back} \\
\prompt{Select filter} & \key{1} maximum resolution, \key{2} method,
  \key{3} contains metal, \key{4} clear (does not work), \key{5} back \\
AlphaFold / AlphaFill & \key{1} UniProt ID, \key{2} gene name, \key{3} back \\
File browser & a number, \typed{..}, a full path, \key{q} (then a path question) \\
1b \prompt{Select option} & a number from REMARK 350, \key{m} manual matrix
  (closes ProPrep in 1.23.0), \key{c} cancel \\
1b \prompt{Include ligands/cofactors (HETATM)?} & \key{y} (default) \\
1b \prompt{Include water molecules?} & \key{n} (default) \\
\bottomrule
\end{longtable}

\section{If something goes wrong}
\label{sec:structure-problems}

\begin{description}[style=nextline]
  \item[``not found (404)''] The full message is \prompt{PDB XQ99 not found
        (404). May be a large structure only available in mmCIF format.}
        Either the code does not exist (check it on the RCSB website) or the
        entry is very large (more than 99\,999 atoms, such as ribosomes and
        viruses), which the RCSB provides only as mmCIF. ProPrep cannot
        download such entries.
  \item[``control characters''] \prompt{URL can't contain control characters}
        means you typed a space or other text that is not a PDB code. Nothing
        is broken; try again.
  \item[``No PDB ID entered''] You pressed \key{Enter} without a code.
  \item[``search service could not be reached''] You have no internet
        connection. A search needs one.
  \item[``No AlphaFold structure found''] The accession may be wrong (then
        the screen also shows \prompt{HTTP Error 400: Bad Request}). Or the
        protein has no prediction, or you are offline.
  \item[``is not a PDB file''] The browser accepts only \file{.pdb} files.
        Convert the file, or download the PDB-format version.
  \item[A tool works on the wrong structure] See
        Section~\ref{sec:structure-switching}. Start a new project.
  \item[``No REMARK 350 records found''] Tool \key{1b} is reading a file
        that has no header lines. Do not choose \key{m}. Instead, download the assembly with
        structure type \key{2} in a new project.
  \item[ProPrep closes in 1b] It closed after \prompt{Enter transformation
        matrix}: the manual matrix does not work in 1.23.0. Restart ProPrep;
        your files are kept (Chapter~\ref{ch:sessions} explains how to resume).
\end{description}

Chapter~\ref{ch:troubleshooting} lists more symptoms, and
Chapter~\ref{ch:knownissues} lists all known limitations of version 1.23.0.
